\section{¿Por qué estudiar el Agent Framework?}

\subsection{De un solo Agent a la colaboración entre múltiples Agents}

\begin{figure}[H]
\centering
\includegraphics[page=6,width=0.9\textwidth]{slides.pdf}
\caption{La introducción describe la colaboración entre múltiples Agents como un «mecanismo de división del trabajo» en tareas complejas, donde distintos módulos se encargan de distintas subtareas.}
\end{figure}

Cuando la duración de la tarea es suficientemente larga y las subtareas son suficientemente distintas entre sí, un solo Agent suele fallar por la saturación del contexto, la mezcla de responsabilidades y la deriva del razonamiento. Por eso el curso coloca la colaboración entre múltiples Agents entre sus temas centrales:
\begin{itemize}
    \item un planner especializado se encarga de descomponer la tarea;
    \item un specialized executor se encarga de operaciones locales como código, recuperación de información y análisis de tablas;
    \item un verifier / reviewer se encarga de revisar la exactitud factual, la consistencia y la seguridad.
\end{itemize}

\begin{importantbox}{El verdadero valor del Framework}
El valor del Agent framework no está en tener varios conjuntos de API, sino en que ofrece una manera estable de organizar tareas complejas: cómo dividir los módulos, cómo transmitir el contexto, cómo hacer que colaboren varios roles y cómo integrar la recuperación de fallos en el flujo de ejecución.
\end{importantbox}

\subsection{Oportunidades de aplicación y presión de evaluación}

\begin{figure}[H]
\centering
\includegraphics[page=8,width=0.9\textwidth]{slides.pdf}
\caption{El curso presenta lado a lado la generación de código, la automatización de flujos de trabajo, los asistentes personales y la robótica, para mostrar que el límite de aplicaciones de los Agents se está expandiendo rápidamente.}
\end{figure}

\begin{figure}[H]
\centering
\includegraphics[page=9,width=0.9\textwidth]{slides.pdf}
\caption{Puntos de referencia como GAIA, WebArena y SWE-Bench constituyen las primeras coordenadas de evaluación de la investigación sobre Agents.}
\end{figure}

El curso presenta las aplicaciones junto con los benchmark desde el principio, porque los sistemas de Agent son fáciles de convertir en un demo, pero difíciles de poner en producción de manera estable. Los retos que enumera la introducción incluyen:
\begin{itemize}
    \item el razonamiento y la planificación son inestables en tareas complejas;
    \item la eficiencia para aprovechar la retroalimentación del entorno es baja en tareas de largo horizonte;
    \item la comprensión multimodal y el world model todavía no son maduros;
    \item la colaboración entre múltiples Agents, la explicación teórica y la seguridad y privacidad todavía no cuentan con un marco sólido.
\end{itemize}

\begin{warningbox}{La advertencia de riesgo más importante de la introducción}
El fallo de un sistema de Agent normalmente no es «responder mal una pregunta», sino un sesgo acumulado a lo largo de una tarea prolongada: un plan equivocado, una llamada incorrecta a una herramienta, una lectura o escritura errónea del entorno externo, que terminan amplificando un problema pequeño hasta convertirlo en una falla sistémica. Por eso desplegar un Agent exige más observabilidad y evaluación que desplegar una aplicación de LLM común.
\end{warningbox}

\subsection{Resumen de la sección}

La razón para estudiar el Agent framework es muy práctica: determina si el sistema puede manejar tareas reales, y también si esos sistemas pueden evaluarse, depurarse y desplegarse de manera estable.
